% =====================================================================
%  theory/characteristic_matrix.tex
%  Phase 1, Blueprint V5.0, Task E
%
%  Construction of M(xi; k, omega) a = 0 for a 6mm piezothermoelastic
%  medium, its exact determinant, the root inventory, parity, gauge roots,
%  admissible-mode selection, and the normal / oblique limits.
%
%  SOURCE-ONLY: no TeX engine is installed. Verified as text.
%  The determinant coefficients quoted below are produced exactly by
%  code/foundation_checks.py (cached in code/det_coeffs.json) and every
%  structural claim is checked numerically in checks C1, C2 and I1.
% =====================================================================
\documentclass[11pt]{article}
\usepackage{amsmath,amssymb,amsthm,bm}
\usepackage{geometry}
\usepackage{booktabs}
\geometry{margin=1in}
\newcommand{\s}{\sigma}
\newcommand{\eps}{\varepsilon}
\newcommand{\bth}{\vartheta}
\newcommand{\D}{\bm{D}}
\newcommand{\E}{\bm{E}}
\newcommand{\bu}{\bm{u}}
\newcommand{\Lam}{\Lambda}
\theoremstyle{plain}
\newtheorem{proposition}{Proposition}
\newtheorem{remark}{Remark}
\newtheorem{corollary}{Corollary}

\title{The characteristic matrix: structure, root inventory and\\
admissible modes\\
\large Theory note 3 of 4 --- Phase 1, Blueprint V5.0}
\author{Derived for the Piezothermoelastic Anti-Reflection programme}
\date{4 November 2026}

\begin{document}
\maketitle

\begin{abstract}
Blueprint V5.0 \S3 states that the number of roots, their parity, the gauge roots
and the mode count are to be \emph{derived, not assumed}, and that a preliminary
toy-coefficient observation of an even polynomial is a prompt rather than
evidence. This note supplies the derivation. The main results are: the
characteristic determinant is a polynomial of degree $8$ in $\xi$ that is
\emph{even} in $\xi$ with exactly five non-zero coefficients, so it factors as
$P_4(\xi^2)$ and the eight roots form four $\pm$ pairs; the four branches are
identified in closed form as quasi-P, quasi-SV, an electrically evanescent branch
and a thermal branch; the only gauge root is the constant-potential state at
$(\kappa,\xi)=(0,0)$; and the count of admissible modes per half-space is $4$,
not $8$. All of this is verified numerically to machine-level residuals by
\texttt{code/foundation\_checks.py}. \textbf{No reflection or transmission
coefficient is computed here.}
\end{abstract}

\section{The ansatz and the matrix}
\label{sec:ansatz}

With the bulk system of \texttt{theory/governing\_equations.tex} and the ansatz
\begin{equation}
  \bm{y} = \bm{a}\,e^{\,i(\kappa x + \xi z - \omega t)},
  \qquad \bm{a} = (U,W,\Phi,\Theta)^{\mathsf T},
  \label{eq:ansatz}
\end{equation}
the four field equations become an algebraic system
\begin{equation}
  \bm{M}(\xi;\kappa,\omega)\,\bm{a} = \bm{0}.
  \label{eq:eigen}
\end{equation}
This is Blueprint \S3. Substituting the kinematics and $E_i=-\varphi_{,i}$
of the governing-equations note into the balance equations gives, entry by entry,
\begin{align}
  M_{11} &= -C_{11}\kappa^2 - C_{44}\xi^2 + \rho\omega^2, \label{eq:M11}\\
  M_{12} &= -\kappa\xi\,(C_{13}+C_{44}), \label{eq:M12}\\
  M_{13} &= -\kappa\xi\,(e_{15}+e_{31}), \label{eq:M13}\\
  M_{14} &= -i\beta_1\kappa, \label{eq:M14}\\[2pt]
  M_{21} &= M_{12}, \label{eq:M21}\\
  M_{22} &= -C_{33}\xi^2 - C_{44}\kappa^2 + \rho\omega^2, \label{eq:M22}\\
  M_{23} &= -e_{15}\kappa^2 - e_{33}\xi^2, \label{eq:M23}\\
  M_{24} &= -i\beta_3\xi, \label{eq:M24}\\[2pt]
  M_{31} &= M_{13}, \label{eq:M31}\\
  M_{32} &= M_{23}, \label{eq:M32}\\
  M_{33} &= \eps_{11}\kappa^2 + \eps_{33}\xi^2, \label{eq:M33}\\
  M_{34} &= i\,p\,\xi, \label{eq:M34}\\[2pt]
  M_{41} &= T_0\beta_1\kappa\omega\,(i\omega\tau_q - 1), \label{eq:M41}\\
  M_{42} &= T_0\beta_3\omega\xi\,(i\omega\tau_q - 1), \label{eq:M42}\\
  M_{43} &= T_0\omega p\,\xi\,(1 - i\omega\tau_q), \label{eq:M43}\\
  M_{44} &= c_E\omega^2\rho\tau_q + i\,c_E\omega\rho
            - \kappa_{11}\kappa^2 - \kappa_{33}\xi^2 . \label{eq:M44}
\end{align}
\textsc{Derived}. Note the identities $M_{21}=M_{12}$, $M_{31}=M_{13}$,
$M_{32}=M_{23}$: the mechanical--mechanical, mechanical--electrical and
electrical--mechanical blocks are symmetric, which is a consequence of the single
potential.

\begin{remark}[Symmetry is only partial]
\label{rem:partial-sym}
$\bm{M}$ is \emph{not} symmetric: $M_{14}\neq M_{41}$ etc. The asymmetry is
exactly the loss of Hermiticity introduced by the first-order time derivatives in
the entropy balance. It is not a numerical artefact and no symmetrisation should
be applied.
\end{remark}

\section{The determinant: degree, parity, closed form}
\label{sec:det}

\begin{proposition}[Structure of $\det\bm{M}$]
\label{prop:det-structure}
For the matrix \eqref{eq:M11}--\eqref{eq:M44},
\begin{equation}
  \det\bm{M}(\xi;\kappa,\omega)
  = P_0 + P_2\,\xi^2 + P_4\,\xi^4 + P_6\,\xi^6 + P_8\,\xi^8 ,
  \label{eq:det}
\end{equation}
i.e.\ the determinant is a polynomial of degree $8$ in $\xi$ with \textbf{only
even powers present}. The five coefficients are non-zero for generic material
data.
\end{proposition}

\begin{proof}[Proof]
\textbf{Corrected in Phase 1B.} The Phase-1 proof asserted the conjugating matrix
$\bm{S}=\mathrm{diag}(1,1,1,-1)$. That is \emph{wrong} and is withdrawn: it does
not satisfy the entrywise conditions, as is shown below.

Classify every entry by its parity in $\xi$:
\begin{equation}
  \text{even: }(1,1),(2,2),(3,3),(4,4),(1,4),(4,1),(2,3),(3,2);
  \qquad
  \text{odd: }(1,2),(2,1),(1,3),(3,1),(2,4),(4,2),(3,4),(4,3).
  \label{eq:parity-classes}
\end{equation}
For $\bm{M}(-\xi)=\bm{S}\bm{M}(\xi)\bm{S}$ with $\bm{S}=\mathrm{diag}(s_1,\dots,s_4)$
to hold entry by entry we need
\begin{equation}
  s_i s_j =
  \begin{cases}
    +1, & (i,j)\ \text{even in }\xi,\\
    -1, & (i,j)\ \text{odd in }\xi,
  \end{cases}
  \qquad\text{for all } 16 \text{ entries.}
  \label{eq:parity-conditions}
\end{equation}
The system \eqref{eq:parity-conditions} has exactly two solutions,
\begin{equation}
  \boxed{\ \bm{S}=\mathrm{diag}(1,-1,-1,1)\ }\quad\text{and its negative,}
  \label{eq:correct-S}
\end{equation}
which are the same symmetry up to an overall sign. Hence
$\bm{M}(-\xi)=\bm{S}\bm{M}(\xi)\bm{S}$ and, since $\det\bm{S}=\pm1$ so that
$\det(\bm{S})^2=1$,
\begin{equation}
  \det\bm{M}(-\xi)=\det(\bm{S})\,\det\bm{M}(\xi)\,\det(\bm{S})
  =\det\bm{M}(\xi).
  \label{eq:det-even}
\end{equation}
The Phase-1 matrix $\mathrm{diag}(1,1,1,-1)$ fails \eqref{eq:parity-conditions} on
six entries --- $(1,2),(1,3),(1,4),(2,1),(3,1),(4,1)$ --- and is not a symmetry of
$\bm{M}$ at all.
\end{proof}

\begin{remark}[Why the Phase-1 error survived a ``check'']
\label{rem:why-missed}
The determinant parity \eqref{eq:det-even} follows from
$\det(\bm{S})^2=1$ for \emph{either} sign pattern, so a determinant-only test
cannot distinguish $\mathrm{diag}(1,1,1,-1)$ from $\mathrm{diag}(1,-1,-1,1)$:
both give an even determinant, and both would ``pass''. The Phase-1 document
claimed that the 16 entries had been evaluated symbolically and found to vanish.
They had not been. The claim was false, and this is the defect the supervisor
identified. The parity of $\det\bm{M}$ was never in doubt; the \emph{reason}
given for it was wrong, and the wrong matrix was quoted as evidence.
\end{remark}

\begin{remark}[Independent confirmation, now actually performed]
\label{rem:parity-check}
Check C1 of \texttt{code/foundation\_checks.py}, rewritten in Phase 1B, builds the
$4\times4$ matrix with \emph{symbolic} material constants (independently of the
numeric material tables and of the cached determinant), evaluates
$\bm{M}(-\xi)-\bm{S}\bm{M}(\xi)\bm{S}$ entry by entry for \emph{both} candidate
matrices, and reports the failures:
\begin{center}
\begin{tabular}{@{}lcc@{}}
\toprule
$\bm{S}$ & entries failing (of 16) & verdict \\
\midrule
$\mathrm{diag}(1,1,1,-1)$ (Phase 1) & 6 & not a symmetry \\
$\mathrm{diag}(1,-1,-1,1)$ (corrected) & 0 & symmetry \\
\bottomrule
\end{tabular}
\end{center}
An exhaustive search over all $2^4$ diagonal sign patterns reproduces
\eqref{eq:correct-S} and nothing else. The determinant parity is then confirmed
separately by fraction-free elimination of $\bm{M}(\xi)-\bm{M}(-\xi)$.
\end{remark}

\begin{remark}[The root inventory is unaffected]
\label{rem:roots-unaffected}
The correction changes a \emph{proof}, not a \emph{result}. The determinant was
already shown to be even by exact expansion of the coefficients
(\eqref{eq:det}), and the eight roots, their branch assignment and their
residuals are unchanged. What changes is that the stated reason is now correct,
and the quoted matrix is now the right one.
\end{remark}

\subsection{Closed forms of the extreme coefficients}

The two extreme coefficients are simple enough to quote and to use as sanity
checks on any implementation:
\begin{align}
  P_8 &= -\,C_{44}\,\kappa_{33}\big(C_{33}\eps_{33} + e_{33}^2\big),
  \label{eq:P8}\\
  P_0 &= \kappa^8\Big[-C_{11}C_{44}\eps_{11}\kappa_{11}
        - C_{11}e_{15}^2\kappa_{11}\Big] + O(\kappa^6).
  \label{eq:P0}
\end{align}
Equation~\eqref{eq:P8} is worth reading physically: the coefficient of the highest
power of $\xi$ factorises into the shear stiffness, the thermal conductivity along
the poling axis, and the \emph{piezoelectrically stiffened} dielectric factor
$C_{33}\eps_{33}+e_{33}^2$. That last factor is the 1-D stiffened modulus familiar
from the thickness-mode transducer problem, and its appearance here is a useful
sign that the coupling has been assembled correctly.

\subsection{Decoupling structure at $\xi=0$ and at $\kappa=0$}

\begin{proposition}[Block decomposition]
\label{prop:blocks}
\begin{enumerate}
  \item At $\xi=0$ (propagation normal to the poling axis, in the plane of
        transverse isotropy), $\bm{M}$ block-decomposes as
        $\{u_x,\bth\}\oplus\{u_z,\varphi\}$, with all cross terms between the two
        blocks vanishing identically. In the $\{u_x,\bth\}$ block there is
        \emph{no} piezoelectric coupling; in the $\{u_z,\varphi\}$ block there is
        \emph{no} thermal coupling.
  \item At $\kappa=0$ (normal incidence), $\bm{M}$ block-decomposes as
        $\{u_x\}\oplus\{u_z,\varphi,\bth\}$: $u_x$ is completely decoupled and is a
        pure shear mode.
\end{enumerate}
\end{proposition}
\begin{proof}
Direct substitution. At $\xi=0$, $M_{12}=M_{13}=M_{14}=0$ (all carry a factor
$\xi$) and $M_{42}=M_{43}=0$ (both carry $\xi$), leaving $M_{11},M_{41},M_{44}$
in the first block and $M_{22},M_{23},M_{32},M_{33}$ in the second. At
$\kappa=0$, $M_{12}=M_{13}=M_{14}=0$ and $M_{41}=0$, leaving $M_{11}$ isolated.
\end{proof}
\textsc{Tested} (check C1).

\begin{remark}[The thermal field does not couple to shear]
\label{rem:sh}
Combining the two limits: $u_x$ decouples from everything at $\kappa=0$, and the
$\{u_x,\bth\}$ block at $\xi=0$ contains no piezoelectric coupling. The
consequence, stated already in the governing-equations note, is that
\textbf{anti-plane SH motion in a $6mm$ medium with poling along $z$ is thermally
decoupled identically}. SH configurations cannot demonstrate thermal coupling.
Any claim to the contrary would be a modelling error.
\end{remark}

\section{Root inventory and admissible modes}
\label{sec:roots}

\subsection{The eight roots}

By Proposition~\ref{prop:det-structure}, $\det\bm{M}$ factors as $P_4(\xi^2)$, so
the eight roots $\xi_1,\dots,\xi_8$ are four pairs $\pm\xi$. Writing
$\xi = \xi' + i\xi''$ and using $e^{i\xi z}$ with $e^{-i\omega t}$:
\begin{itemize}
  \item $\xi''>0$: the mode decays as $z\to+\infty$ and grows as $z\to-\infty$;
  \item $\xi''<0$: the opposite;
  \item $\xi''=0$: the mode propagates, with $\xi'>0$ travelling in $+z$ and
        $\xi'<0$ in $-z$.
\end{itemize}
Therefore \textbf{exactly four of the eight roots are admissible on each side of an
interface}: those with $\xi''>0$ (or $\xi'>0$ when $\xi''=0$) for medium 2, and
their negatives for medium 1. The count of admissible modes per half-space is
$4$, not $8$. This is the derivation that Blueprint \S3 asks for.

\subsection{Closed-form identification of the four branches}

At a representative point ($f=1$ MHz, $\theta=30^\circ$, PZT-5H) the eight roots
group into four branches, each identified against an independent closed form:

\begin{center}
\begin{tabular}{llll}
\toprule
Branch & Closed form for $\xi$ & Value at the test point & $\lvert\bm{a}\rvert$ pattern \\
\midrule
quasi-P & $\sqrt{(\omega/c_L)^2-\kappa^2}$,
  $c_L^2=(C_{33}+e_{33}^2/\eps_{33})/\rho$ & $1.139\times10^{3}$ & $(0,0,1,\cdot)$ \\
quasi-SV & $\sqrt{(\omega/c_T)^2-\kappa^2}$, $c_T^2=C_{44}/\rho$
  & $3.498\times10^{3}$ & $(0,1,0,0)$ \\
electrically evanescent & $i\kappa\sqrt{\eps_{11}/\eps_{33}}$
  & $i\,8.452\times10^{2}$ & $(0,0,1,0)$ \\
thermal (LS) & $(1+i)\sqrt{\omega\rho c_E/(2\kappa_{33})}$
  & $(-2.569+2.569i)\times10^{6}$ & $(\cdot,\cdot,\cdot,1)$ \\
\bottomrule
\end{tabular}
\end{center}

\textsc{Tested} (check C2): each root is matched to its closed form, and the
eigenmode residual $\lVert\bm{M}\bm{a}\rVert/(\lVert\bm{M}\rVert\lvert\bm{a}\rvert)$
is below $10^{-26}$ for all eight roots.

\begin{remark}[The ``quasi'' in quasi-P and quasi-SV]
\label{rem:quasi}
The closed forms above neglect the piezoelectric and thermal corrections to the
wave speeds; the actual roots differ from them by a few parts in $10^{3}$
(measured: $1.9\%$ for quasi-P, $0.6\%$ for quasi-SV at the test point). They are
labelled ``quasi'' for this reason, and the identification is done by nearest
match, not by assuming equality. A method that assumes the closed forms are exact
would mis-assign the reflected branches near a degeneracy.
\end{remark}

\subsection{The thermal branch is always strongly evanescent}

At $1$ MHz and $T_0=300$ K the thermal wavenumber is
$\lvert\xi_{\mathrm{th}}\rvert \approx 2.6\times10^{6}\ \mathrm{m^{-1}}$ against
$\kappa\approx 7.9\times10^{2}\ \mathrm{m^{-1}}$, i.e.\ the thermal mode decays
over a diffusion length
\begin{equation}
  \mu = \frac{1}{\operatorname{Re}\xi_{\mathrm{th}}}
  = \sqrt{\frac{2\kappa_{33}}{\omega\rho c_E}} \approx 3.9\times10^{-7}\ \mathrm{m}
  \label{eq:mu}
\end{equation}
at $1$ MHz. \textbf{Consequence for the scattering problem:} an evaluation plane
placed more than a few $\mu$ from the interface sees no thermal contribution at
all. Any energy balance must therefore either place its evaluation planes within
about a micrometre of the interface (impractical at these wavelengths) or
demonstrate that the thermal mode contributes a negligible flux at the chosen
planes. The latter is the only practical route at megahertz frequencies, and it
must be demonstrated rather than assumed. This is recorded as a Phase-2
requirement.

\section{Gauge and zero roots}
\label{sec:gauge}

\begin{proposition}[The only gauge root]
\label{prop:gauge}
At $(\kappa,\xi)=(0,0)$ the null space of $\bm{M}$ is one-dimensional, spanned by
$\bm{a}_0=(0,0,1,0)^{\mathsf T}$ --- the constant potential with $D_z=0$. This
state carries no field energy and no flux. It does not persist for $\kappa\neq0$.
\end{proposition}
\begin{proof}
Substituting $\kappa=\xi=0$ into \eqref{eq:M11}--\eqref{eq:M44} leaves a diagonal
matrix with entries $(0,\ 0,\ 0,\ c_E\omega^2\rho\tau_q + ic_E\omega\rho)$, of
which the third vanishes, giving a single null direction. For $\kappa\neq0$ the
entry $M_{33}=\eps_{11}\kappa^2\neq0$, so the null direction is removed.
\end{proof}
\textsc{Tested} (check C1 computes the null space symbolically and obtains exactly
that vector).

\begin{remark}[Two further degeneracies at $\kappa=0$ when $e=0$]
\label{rem:k0-degenerate}
With $\kappa=0$ \emph{and} the piezoelectric coupling set to zero, the coefficients
$P_0$ and $P_2$ of \eqref{eq:det} acquire a factor $\kappa^2$ (visible in
\eqref{eq:P0}), so $\xi=0$ becomes a root of multiplicity four rather than one.
This is not a different gauge root --- it is the same constant-potential
degeneracy, now fourfold. Practically it means: \textbf{normal incidence in the
decoupled limit is a singular point of the eigenvalue problem and must be reached
as a limit} $\kappa\to0$ with the coupling kept non-zero. This was found
numerically (see U-3 in \texttt{reports/UNRESOLVED\_ISSUES.md}) after an attempt
to validate at exactly $\kappa=0$ failed.
\end{remark}

\section{Selection of admissible and outgoing modes}
\label{sec:selection}

Blueprint \S3 requires selection by flux and decay, \emph{not} by the sign of
$\operatorname{Re}\xi$ alone. The rule is:
\begin{enumerate}
  \item \textbf{Decay:} on the medium occupying $z>0$, keep roots with
        $\xi''>0$; on $z<0$, keep $\xi''<0$.
  \item \textbf{Propagating roots ($\xi''=0$):} keep the outgoing direction,
        $\xi'>0$ for $z>0$ and $\xi'<0$ for $z<0$.
  \item \textbf{Verification:} for each kept mode, the cycle-averaged $z$-flux
        $\langle P_z\rangle$ (see \texttt{theory/energy\_flux\_derivation.tex})
        must be positive for a mode outgoing into $z>0$ and negative for one
        outgoing into $z<0$. A root that passes the decay test but fails the flux
        test is a sign error or a wrong branch, not a physical mode.
\end{enumerate}
\textsc{Derived}, and \textsc{Tested} through the antisymmetry
$\langle P_z\rangle(\xi) = -\langle P_z\rangle(-\xi)$, which check D1 verifies
with a maximum relative deviation of $3.5\times10^{-13}$ over all branches of
three materials.

\begin{remark}[Why $\operatorname{Re}\xi$ alone is not enough]
\label{rem:not-re-enough}
For a lossy medium a root can have $\operatorname{Re}\xi>0$ and still be
incoming, because the stored-energy flux and the group velocity can differ in
sign near a strong absorption line. The Lord--Shulman thermal branch is exactly
such a case: it is evanescent, yet the sign of $\operatorname{Re}\xi$ flips
between the two members of the pair while the flux keeps its physical direction.
Using $\operatorname{Re}\xi$ alone would label a decaying mode as incoming.
\end{remark}

\section{Normal and oblique limits}
\label{sec:limits}

\paragraph{Oblique incidence ($\kappa\neq0$).} All four branches carry all four
field components to some degree, and the mode-conversion structure is genuine:
an incident quasi-P wave generates reflected quasi-P \emph{and} quasi-SV, plus
evanescent electrical and thermal near-fields. The number of reflected propagating
channels is $2$, not $3$, because the electrical and thermal branches are
evanescent at these frequencies.

\paragraph{Normal incidence ($\kappa\to0$).} By
Proposition~\ref{prop:blocks}(2), $u_x$ decouples exactly and the remaining
$3\times3$ system $\{u_z,\varphi,\bth\}$ governs the longitudinal problem. This is
the limit in which the thermal and electrical effects are strongest, because the
longitudinal branch is the only one that couples to $\bth$ and $\varphi$.

\paragraph{Why $\kappa=0$ must be a limit.} Proposition~\ref{prop:gauge} and
Remark~\ref{rem:k0-degenerate} show that $\kappa=0$ is a degenerate point of the
eigenvalue problem. Practically: compute at $\kappa=\omega\sin\theta/c$ for
$\theta=10^{-4},10^{-3},10^{-2}$ rad, verify that the quantity of interest
converges as $\theta^2$, and report the extrapolated value. Check I1 does exactly
this against the closed-form elastic branches and obtains agreement to
$1.7\times10^{-15}$.

\section{What this note does and does not establish}
\label{sec:summary}

\paragraph{Established.} The matrix \eqref{eq:M11}--\eqref{eq:M44}; the
determinant structure \eqref{eq:det}, including exact parity with a
proof; the closed forms \eqref{eq:P8}--\eqref{eq:P0}; the block decompositions;
the root count $8$ with $4$ admissible per half-space; the closed-form
identification of all four branches with residuals below $10^{-26}$; the single
gauge root; the selection rule and its verification.

\paragraph{Not established.}
\begin{itemize}
  \item No scattering coefficient. Root structure is not a scattering result.
  \item The branch identification is by nearest match to a closed form that
        neglects coupling corrections of order $10^{-3}$; near a degeneracy of two
        branches this identification would need the coupled dispersion relation
        rather than the closed form.
  \item The thermal-mode contribution to the flux at practical evaluation planes
        is not quantified here; that is a Phase-2 requirement
        (\S\ref{sec:roots}, \eqref{eq:mu}).
  \item The 6 mm restriction is essential. A lower symmetry (e.g.\ $mm2$ or
        triclinic) would break the parity argument \eqref{eq:parity-identity} and
        the block structure, and the degree would remain $8$ but the roots would
        no longer come in $\pm$ pairs in the same way.
\end{itemize}

\end{document}
